\documentclass[10pt,a4paper]{amsart}
\usepackage[margin=19mm]{geometry}
\usepackage{amsmath,amssymb,mathtools}
\usepackage[scr=rsfs,cal=euler]{mathalfa}
\usepackage{xfrac}
\usepackage[unicode,hidelinks,bookmarks=false]{hyperref}
\newcommand{\ie}{i.e.\ }
\newcommand{\cf}{cf.\ }
\newcommand{\resp}{respectively\ }
\newcommand{\Subsection}{\subsection}
\providecommand{\nobreakdash}{\nobreak\hbox{-}}
\setlength{\emergencystretch}{2em}
\multlinegap0pt
\usepackage{amssymb}
\usepackage{mathtools}
\usepackage{bm}
\usepackage[shortlabels]{enumitem}
\newtheorem{lemma}{Lemma}
\newcommand{\R}{\mathbb{R}}
\newcommand{\dd}{\,\mathrm{d}}
\title{A Quadratic-Form Representation of the Scalar Casimir Trace from Codimension-Three Riesz Reduction}
\author{Irshadullah Khan, Bilal Khan}
\date{Source: arXiv:2605.06693v2 (CC BY 4.0)}
\begin{document}
\maketitle

\noindent\emph{Source and licence.} A Quadratic-Form Representation of the Scalar Casimir Trace from Codimension-Three Riesz Reduction. Version arXiv:2605.06693v2 (CC BY 4.0), distributed by arXiv under the Creative Commons Attribution 4.0 International licence, http://creativecommons.org/licenses/by/4.0/, as recorded in the arXiv licence metadata for this exact version and shown on the abstract page https://arxiv.org/abs/2605.06693v2. Full licence notice: \texttt{licenses/arxiv-2605.06693-CC-BY-4.0.txt}.\par\smallskip

\noindent\textit{Editorial scope.} Counted unit: the article's lemma lem:gaussian-interval-log-concavity (source0.tex lines 2317--2324). The article's complete original proof of it is delivered with it (source0.tex lines 2326--2445, one \texttt{proof} environment of 120 lines). No mathematics is written, repaired or re-derived: the statement and the proof are the article's own lines, in the article's own notation, and no retained line is substituted or re-flowed.  No further source-local context is needed: every cross-reference of every retained line resolves inside the two delivered spans.  Nothing was omitted from any retained span, so the package carries no omission record.  No retained line contains a citation, so the wrapper carries no bibliography and no bibliography entry.  No retained line contains a figure environment, an image-inclusion command or drawing code, so no image is delivered and the package declares no asset dependency.  Editorial formatting only, in addition: an AMS \texttt{amsart} preamble that restates the article's own declarations and macros used by the retained lines, namely the environments \texttt{lemma} on separate counters, together with the standard packages those lines need.  The retained environments print in this document as Lemma 1 in order of appearance, whereas the article prints the same items with the numbering of its own full document.  The complete native source of the article as distributed in the arXiv e-print for this exact version is bundled unchanged as \texttt{sources/200042/source0.tex}.
\begin{lemma}[Strict log-concavity in logarithmic length]
\label{lem:gaussian-interval-log-concavity}
For every fixed \(t>0\), the function
\[
    u\longmapsto \log I_{e^u}(t)
\]
is strictly concave on \(\R\).
\end{lemma}

\begin{proof}
By the change of variables \(X=\sqrt t\,x\), \(Y=\sqrt t\,y\),
\[
    I_L(t)=t^{-1}J(L\sqrt t),
\]
where
\[
    J(r):=\int_0^r\int_0^r e^{-(X-Y)^2}\,\dd X\,\dd Y
    =2\int_0^r(r-s)e^{-s^2}\,\dd s.
\]
Multiplying by the positive constant \(t^{-1}\) and translating the variable
\(u\) by \((1/2)\log t\) do not affect strict concavity.  It is therefore
enough to prove that
\[
    w\longmapsto \log J(e^w)
\]
is strictly concave.

Since the factor \(2\) in \(J\) is irrelevant for logarithmic concavity, set
\[
    j(r):=\frac{1}{2}J(r)=\int_0^r(r-s)e^{-s^2}\,\dd s.
\]
Introduce
\[
    A(r):=\int_0^r e^{-s^2}\,\dd s,
    \qquad
    E(r):=e^{-r^2},
    \qquad
    B(r):=\int_0^r s e^{-s^2}\,\dd s=\frac{1-E(r)}{2}.
\]
Then
\[
    j(r)=rA(r)-B(r),
    \qquad
    j'(r)=A(r),
    \qquad
    j''(r)=E(r).
\]
Let \(r=e^w\).  A direct differentiation gives
\begin{equation}
    \frac{\dd^2}{\dd w^2}\log j(e^w)
    =
    \frac{r(A(r)+rE(r))j(r)-r^2A(r)^2}{j(r)^2}.
    \label{eq:appendix-H-second-derivative}
\end{equation}
Thus it remains to prove that the numerator in
\eqref{eq:appendix-H-second-derivative} is negative for every \(r>0\).
Equivalently, since \(j(r)=rA(r)-B(r)\), we must prove
\begin{equation}
    B(r)(A(r)+rE(r))-r^2A(r)E(r)>0.
    \label{eq:appendix-core-positivity}
\end{equation}
Define
\[
    h(r):=(1-E(r))(A(r)+rE(r))-2r^2A(r)E(r).
\]
Then \eqref{eq:appendix-core-positivity} is equivalent to \(h(r)>0\).  Since
\(h(0)=0\), it is enough to show \(h'(r)>0\) for \(r>0\).  Differentiating,
using \(A'(r)=E(r)\) and \(E'(r)=-2rE(r)\), gives
\begin{equation}
    h'(r)=2E(r)k(r),
    \label{eq:appendix-h-prime}
\end{equation}
where
\begin{equation}
    k(r):=rA(r)(2r^2-1)+(1-r^2)(1-E(r)).
    \label{eq:appendix-k-definition}
\end{equation}
We prove \(k(r)>0\) in two ranges.

First let \(0<r\leq 1/\sqrt2\).  Then \(2r^2-1\leq0\), and
\(A(r)\leq r\), so
\[
    rA(r)(2r^2-1)
    \geq
    r^2(2r^2-1).
\]
Also, with \(x=r^2\), the elementary inequality
\(1-e^{-x}\geq x-x^2/2\) gives
\[
    1-E(r)\geq r^2-\frac{r^4}{2}.
\]
Therefore
\[
\begin{split}
    k(r)
    &\geq
    r^2(2r^2-1)
    +(1-r^2)\left(r^2-\frac{r^4}{2}\right)  \\
    &=
    \frac12 r^4(1+r^2)>0.
\end{split}
\]

Now let \(r\geq1/\sqrt2\).  Then \(2r^2-1\geq0\).  Since
\[
    1-E(r)
    =\int_0^r 2s e^{-s^2}\,\dd s
    \leq
    2r\int_0^r e^{-s^2}\,\dd s
    =2rA(r),
\]
we have \(A(r)\geq(1-E(r))/(2r)\).  Hence
\[
\begin{split}
    k(r)
    &\geq
    \frac{1-E(r)}{2}(2r^2-1)
    +(1-r^2)(1-E(r))  \\
    &=
    \frac{1-E(r)}{2}>0.
\end{split}
\]
Thus \(k(r)>0\) for all \(r>0\).  By \eqref{eq:appendix-h-prime},
\(h'(r)>0\) for all \(r>0\), and since \(h(0)=0\), \(h(r)>0\) for all
\(r>0\).  This proves \eqref{eq:appendix-core-positivity}, hence the second
derivative in \eqref{eq:appendix-H-second-derivative} is strictly negative.
Therefore \(w\mapsto\log J(e^w)\), and hence
\(u\mapsto\log I_{e^u}(t)\), is strictly concave.
\end{proof}

\end{document}
